\chapter*{\Large \center Abstract}

Prompt-based large language models (LLMs), such as GPT-style systems, Gemini, and LLaMA, demonstrate strong zero-shot and few-shot performance in Grammatical Error Correction (GEC). However, their generative fluency can lead to over-correction, whereby models introduce unnecessary lexical or structural changes beyond minimal grammatical correction. For second-language (L2) learners, such rewriting may obscure the source of errors and reduce the pedagogical value of feedback~\cite{loem2023exploringeffectivenessgpt3grammatical}.

This dissertation investigates how prompt engineering influences over-correction risk in LLM-based GEC. It compares multiple prompting strategies and evaluates model outputs using ERRANT-based accuracy metrics alongside a Composite Over-Correction Index (OCI), which is used as a comparative indicator of source-output rewriting after accounting for estimated utility.

The findings suggest that prompt design plays an important role in how LLMs correct learner writing. Within the evaluated setup, instruction-constrained and role-based prompts reduced OCI while maintaining strong grammatical accuracy. Overall, the study provides a structured way to analyse over-correction risk and offers guidance for designing LLM-based writing support tools that make minimal but useful corrections. The broader aim is to support L2 learners' writing development without replacing their own voice~\cite{park-etal-2025-leveraging}.
